\paragrafos{Ευκλείδεια διαίρεση: το υπόλοιπο που μένει}{Α.1.2, «Ταυτότητα Ευκλείδειας Διαίρεσης»}

\begin{anagke}
Το σχολείο οργανώνει την πρώτη του εκδρομή. Θα πάνε και τα 387 παιδιά, με λεωφορεία των 49 θέσεων. Η κυρία Αλεξάνδρα ζητάει από την τάξη να βρει πόσα λεωφορεία χρειάζονται.

Ο Στέλιος, που αγαπάει το κομπιουτεράκι, κάνει $387:49$ και η οθόνη γράφει $7{,}897959\dots$ «Εφτά κόμμα εννιά!» λέει περήφανος. «Άρα εφτά λεωφορεία». Η Ρήνη τον κοιτάει. «Και τα υπόλοιπα παιδιά; Μένουν στην αυλή;» Ο Στέλιος ξαναδιαβάζει την οθόνη. «Μα το 0,897 τι είναι; Δεν υπάρχει 0,897 λεωφορείου». Η Ζωή κάνει τη διαίρεση στο χαρτί, όπως τη μάθαμε: πηλίκο 7 και κάτι περισσεύει. «Αυτό που περισσεύει είναι παιδιά, όχι κόμματα. Εκεί βρίσκεται η απάντηση».

Σε πολλά προβλήματα δεν ενδιαφέρει μόνο πόσες φορές χωράει ο ένας αριθμός στον άλλον, ενδιαφέρει και τι περισσεύει. Σε αυτή την παράγραφο μαθαίνουμε να γράφουμε τη διαίρεση με μία ισότητα, την ταυτότητα της Ευκλείδειας διαίρεσης, και να τη χρησιμοποιούμε για να λύνουμε προβλήματα.
\end{anagke}

\begin{lexi}{πηλίκο, υπόλοιπο, Ευκλείδεια}{πηλίκο, υπόλοιπο}
Το \lexiword{πηλίκο} έρχεται από το αρχαίο «πηλίκος», που σημαίνει «πόσο μεγάλος»: η διαίρεση απαντά στο «πόσο μεγάλο μερίδιο παίρνει ο καθένας». Το \lexiword{υπόλοιπο} είναι κυριολεκτικά αυτό που μένει, το «λοιπό» που δεν μοιράστηκε.

Η \lexiword{Ευκλείδεια διαίρεση} πήρε το όνομά της από τον Ευκλείδη, τον Έλληνα μαθηματικό που έζησε στην Αλεξάνδρεια γύρω στο 300 π.Χ. και έγραψε τα «Στοιχεία», το βιβλίο όπου συγκέντρωσε τα μαθηματικά της εποχής του.
\end{lexi}

\begin{metafrasi}
Μοιράζω 43 μπάλες σε ομάδες των 7: γίνονται 6 ομάδες και μένει 1 μπάλα. & $43=7\cdot 6+1$\\
Δεν περισσεύει τίποτα (τέλεια διαίρεση). & $\Delta=\delta\cdot\pi$, υπόλοιπο $0$\\
Το υπόλοιπο είναι πάντα μικρότερο από τον διαιρέτη. & $\upsilon<\delta$\\
Στη διαίρεση με το 8 τα πιθανά υπόλοιπα. & $0,1,2,3,4,5,6,7$\\
Δεν διαιρούμε ποτέ με το μηδέν. & $\delta\neq 0$\\
\end{metafrasi}

\tmima{Η θεωρία}

\begin{kanonas}{ταυτότητα της Ευκλείδειας διαίρεσης}
Αν $\Delta$ και $\delta$ είναι δύο φυσικοί αριθμοί με $\delta\neq 0$, τότε υπάρχουν δύο άλλοι φυσικοί $\pi$ και $\upsilon$ έτσι ώστε
$$\Delta=\delta\cdot\pi+\upsilon,\qquad\text{με}\quad \upsilon<\delta.$$
\end{kanonas}

Ο $\Delta$ λέγεται \textbf{διαιρετέος}, ο $\delta$ \textbf{διαιρέτης}, ο $\pi$ \textbf{πηλίκο} και ο $\upsilon$ \textbf{υπόλοιπο}. Για παράδειγμα, από τη διαίρεση $43:7$ έχουμε $\pi=6$ και $\upsilon=1$, δηλαδή $43=7\cdot 6+1$ με $1<7$.

Αν το υπόλοιπο είναι $0$, η διαίρεση λέγεται \textbf{τέλεια} και γράφεται $\Delta=\delta\cdot\pi$. Τότε ισχύει και $\Delta:\delta=\pi$ και $\Delta:\pi=\delta$. Η τέλεια διαίρεση είναι η πράξη αντίστροφη του πολλαπλασιασμού.

\begin{idiotita}{ειδικές περιπτώσεις}
\begin{itemize}[leftmargin=1.2em,itemsep=0.3ex,topsep=2pt]
\item Ο διαιρέτης δεν μπορεί να είναι $0$: $\delta\neq 0$.
\item $\alpha:\alpha=1$ (όταν $\Delta=\delta$ το πηλίκο είναι $1$).
\item $\alpha:1=\alpha$ (όταν $\delta=1$ το πηλίκο είναι ο διαιρετέος).
\item $0:\alpha=0$ (όταν $\Delta=0$ το πηλίκο είναι $0$).
\end{itemize}
\end{idiotita}

\begin{prosoxi}
Μια ισότητα της μορφής $\Delta=\delta\cdot\pi+\upsilon$ δεν είναι πάντα Ευκλείδεια διαίρεση. Πρέπει να ισχύει και $\upsilon<\delta$. Για παράδειγμα η $50=7\cdot 6+8$ είναι αληθινή ισότητα, αλλά δεν είναι Ευκλείδεια διαίρεση, γιατί $8>7$: στο $50$ το $7$ χωράει ακόμη μία φορά ($50=7\cdot 7+1$).
\end{prosoxi}

\tmimasmall{Τι υπόλοιπα μπορούν να βγουν}

Επειδή $\upsilon<\delta$, στη διαίρεση με το $5$ τα πιθανά υπόλοιπα είναι μόνο $0,1,2,3,4$. Στη διαίρεση με το $\delta$ υπάρχουν ακριβώς $\delta$ διαφορετικά υπόλοιπα, από το $0$ ως το $\delta-1$. Γι' αυτό τα προβλήματα που επαναλαμβάνονται κυκλικά (ημέρες της εβδομάδας, ώρες του ρολογιού) λύνονται με Ευκλείδεια διαίρεση.

\tmima{Λυμένα παραδείγματα}

\paradeigma{Ποιες ισότητες είναι Ευκλείδεια διαίρεση;}
\tixero Οι ισότητες α) $54=25\cdot 3+9$, β) $370=46\cdot 8+2$, γ) $680=51\cdot 12+68$.\par
\tithelo Σε ποιες ισχύει $\upsilon<\delta$ και ποιοι είναι οι $\Delta,\delta,\pi,\upsilon$;\par
\possindeo Στην α) είναι $9<25$: Ευκλείδεια διαίρεση με $\Delta=54$, $\delta=25$, $\pi=3$, $\upsilon=9$. Στη β) το $2$ είναι μικρότερο και από τον $46$ και από τον $8$, άρα μπορούμε να πούμε $\delta=46,\pi=8$ ή $\delta=8,\pi=46$ (και στις δύο $\upsilon=2$). Στη γ) το $68$ είναι μεγαλύτερο και από τον $51$ και από τον $12$, άρα δεν είναι διαίρεση.\par
\apantisi Οι α) και β) εκφράζουν Ευκλείδεια διαίρεση, η γ) όχι.

\paradeigma{Οι καρέκλες της αίθουσας}
\tixero Κάθε αίθουσα πρέπει να έχει $12$ θρανία. Σήμερα ήρθαν $3$ φορτηγά που το καθένα μετέφερε $35$ θρανία.\par
\tithelo Πόσες αίθουσες συμπληρώνονται και πόσα θρανία χρειάζονται ακόμη για να συμπληρωθεί μία ακόμη;\par
\possindeo Συνολικά $3\cdot 35=105$ θρανία. Κάνω τη διαίρεση $105:12$: $105=12\cdot 8+9$. Άρα $8$ αίθουσες συμπληρώνονται και περισσεύουν $9$ θρανία. Για την επόμενη χρειάζονται ακόμη $12-9=3$.\par
\apantisi $8$ αίθουσες, και χρειάζονται άλλα $3$ θρανία.\par
\elegxos $12\cdot 8+9=96+9=105$ και $9<12$.

\paradeigma{Η μέρα της εβδομάδας}
\tixero Σήμερα είναι Δευτέρα.\par
\tithelo Τι μέρα θα είναι σε $100$ ημέρες;\par
\possindeo Οι ημέρες επαναλαμβάνονται ανά $7$. Κάνω τη διαίρεση: $100=7\cdot 14+2$. Μετά από $14$ πλήρεις εβδομάδες ξαναγυρνάμε σε Δευτέρα, και μένουν $2$ ημέρες: Τρίτη, Τετάρτη.\par
\apantisi Τετάρτη.\par
\elegxos $7\cdot 14=98$ ημέρες είναι Δευτέρα, και $2$ ακόμη φέρνουν Τετάρτη.

\tmima{Ασκήσεις}

\zoni{Το κατάλαβα}{εφαρμογή της ταυτότητας}

\ex Συμπλήρωσε: α) Στη διαίρεση $47:6$ ο διαιρετέος είναι \dots, ο διαιρέτης \dots, το πηλίκο \dots και το υπόλοιπο \dots β) Ο διαιρέτης δεν μπορεί να είναι \dots γ) Αν ο διαιρετέος ισούται με τον διαιρέτη, το πηλίκο είναι \dots δ) Αν ο διαιρετέος είναι $0$, το πηλίκο είναι \dots

\ex Ποιες από τις ισότητες εκφράζουν Ευκλείδεια διαίρεση; α) $47=6\cdot 7+5$ \quad β) $50=7\cdot 6+8$ \quad γ) $29=4\cdot 7+1$ \quad δ) $100=9\cdot 10+10$ \quad ε) $63=9\cdot 7+0$.

\ex Κάνε τις διαιρέσεις και τις δοκιμές τους: α) $135:5$ \quad β) $95:4$ \quad γ) $276:12$ \quad δ) $1.000:7$. Ποιες είναι τέλειες;

\ex Ποια υπόλοιπα μπορεί να έχει μια διαίρεση με διαιρέτη το $5$; Ποια με διαιρέτη το $8$; Ποια με διαιρέτη το $1$;

\ex Σωστό ή λάθος; α) Η διαίρεση $452:2$ είναι τέλεια. \quad β) Ο διαιρέτης μπορεί να είναι $0$. \quad γ) Το υπόλοιπο είναι πάντα μικρότερο από τον διαιρέτη. \quad δ) Αν ο διαιρέτης είναι $1$, το πηλίκο ισούται με τον διαιρετέο. \quad ε) $0:7=7$. \quad στ) $\alpha:\alpha=\alpha$.

\ex Γράψε την ταυτότητα της Ευκλείδειας διαίρεσης για: α) $83:9$ \quad β) $250:25$ \quad γ) $7:10$ \quad δ) $0:5$.

\zoni{Το χρησιμοποιώ}{προβλήματα από τη ζωή}

\ex \textbf{Τα βαρέλια.} Ο παππούς Γιάννης έχει $350$ λίτρα κρασί και βάζει $12$ λίτρα σε κάθε βαρέλι. Πόσα βαρέλια γεμίζουν και πόσα λίτρα περισσεύουν;

\ex \textbf{Η εκδρομή.} Στην εκδρομή θέλουν να πάνε $250$ παιδιά και κάθε λεωφορείο έχει $45$ θέσεις. Πόσα λεωφορεία χρειάζονται; Πόσες θέσεις μένουν άδειες στο τελευταίο;

\ex \textbf{Οι ημέρες.} Σήμερα είναι Τετάρτη. Τι μέρα θα είναι μετά από $45$ ημέρες; Και μετά από $100$ ημέρες;

\ex \textbf{Τα βραχιόλια.} Η Ρήνη έχει $100$ χάντρες και φτιάχνει βραχιόλια με $8$ χάντρες το καθένα. Πόσα βραχιόλια φτιάχνει και πόσες χάντρες περισσεύουν; Πόσες χάντρες της λείπουν για να φτιάξει ένα ακόμη;

\ex \textbf{Το όνομα που επαναλαμβάνεται.} Η Αγγελική γράφει συνέχεια τη λέξη $\mathrm{ΑΓΓΕΛΙΚΗ}$ σε μια λωρίδα χαρτί: $\mathrm{ΑΓΓΕΛΙΚΗΑΓΓΕΛΙΚΗ}\dots$ Ποιο γράμμα βρίσκεται στην $100$ή θέση;

\zoni{Το σκέφτομαι}{ερωτήσεις με λίγη περισσότερη σκέψη}

\ex Ένας αριθμός διαιρούμενος με το $7$ δίνει πηλίκο $12$ και υπόλοιπο $6$. Ποιος είναι ο αριθμός;

\ex Ένας αριθμός διαιρείται τέλεια με το $14$ και δίνει πηλίκο $201$. Ποιος είναι ο αριθμός;

\ex Ένας αριθμός διαιρούμενος με το $9$ έχει πηλίκο $5$ και το μεγαλύτερο δυνατό υπόλοιπο. Ποιος είναι; Τι πηλίκο και τι υπόλοιπο δίνει ο επόμενός του;

\ex Ο Βασίλης λέει: «Αν ένας αριθμός διαιρούμενος με το $4$ δίνει υπόλοιπο $3$, τότε ο διπλάσιός του διαιρούμενος με το $4$ θα δώσει υπόλοιπο $6$». Εξήγησε γιατί αυτό δεν μπορεί να ισχύει και βρες το σωστό υπόλοιπο. (Δοκίμασε με τον αριθμό $7$.)

\zoni{Για το σπίτι}{εξάσκηση}

\ex Υπολόγισε τις τιμές: α) $35-12:3+3\cdot 3$ \quad β) $63:9+63:7-12:2$ \quad γ) $(5+12:3):3-24:(2\cdot 2+16:4)$.

\ex \textbf{Οι συνοδοί.} Στο καταφύγιο ζώων της περιοχής πρέπει να υπάρχει ένας εθελοντής ανά $15$ σκυλάκια, το πολύ. Αν στο καταφύγιο υπάρχουν $162$ σκυλάκια, ποιος είναι ο μικρότερος αριθμός εθελοντών που χρειάζεται;
